% Standalone-includable figure: programs P_1, P_2, P_3 rewritten via the
% abstraction f_0 = \lambda x y. (+ x (* y y)), plus a non-matching example P_3'.
%
% To \input this file, the enclosing document must load:
%   \usepackage{tikz}
%   \usepackage{amsmath}
%   \usepackage[normalem]{ulem}
%   \usepackage{xcolor}
%   \usetikzlibrary{arrows.meta, matrix, calc, backgrounds, fit}

% Darkened + more-saturated versions of the theme pastels #80CDFF and #BD80FF.
\definecolor{myblue}{HTML}{2A7FBF}
\definecolor{mypurple}{HTML}{7230C2}

% Push ulem underlines / wavy lines a bit further below the text baseline so
% they don't crowd descenders or parentheses.
\setlength{\ULdepth}{1.2ex}

% Blue single underline at the normal underline position; text stays black.
\providecommand{\rul}[1]{\textcolor{myblue}{\uline{\textcolor{black}{#1}}}}
% Typeset argument normally and draw a blue rule BELOW the normal underline
% position. Meant to wrap an outer span that contains one or more inner \rul
% regions, so the outer span's line sits deeper than the inner one's.
\ifcsname rullowwd\endcsname\else\newlength{\rullowwd}\fi
\providecommand{\rullow}[1]{%
  \settowidth{\rullowwd}{#1}%
  {#1}%
  \hspace*{-\rullowwd}%
  \textcolor{myblue}{\rule[-2.2ex]{\rullowwd}{0.4pt}}%
}
% Blue wavy underline; text stays black.
\providecommand{\rulw}[1]{\textcolor{myblue}{\uwave{\textcolor{black}{#1}}}}
% Blue operator, used inside underlined regions to color + and * blue.
\providecommand{\bop}[1]{\textcolor{myblue}{#1}}
% Blue math-mode f_0 callsite marker used in the rewritten programs.
\providecommand{\fz}{\textcolor{myblue}{$f_0$}}

\begin{tikzpicture}[>=Stealth]

  \matrix (m) [matrix of nodes,
               row sep=2mm, column sep=72mm,
               nodes={anchor=west, inner sep=3pt}]
  {
    {$P_1$: \texttt{\rul{(\bop{+} u (\bop{*} (/ v 6) (/ v 6)))}}}
      & {$P_1 \mid f_0$: \texttt{(\fz{} u (/ v 6))}} \\
    {$P_2$: \texttt{\rullow{(\bop{+} \rul{(\bop{+} v (\bop{*} v v))} (\bop{*} u u))}}}
      & {$P_2 \mid f_0$: \texttt{(\fz{} (\fz{} v v) u)}} \\
    {$P_3$: \texttt{(* 7 \rul{(\bop{+} 0 (\bop{*} (- v) (- v)))})}}
      & {$P_3 \mid f_0$: \texttt{(* 7 (\fz{} 0 (- v)))}} \\[7mm]
    {$P_3'$: \texttt{(* 7 \rulw{(* (- v) (- v))})}}
      & {$P_3' \mid f_0$: \texttt{(* 7 (* (- v) (- v)))}} \\
  };

  % Coloured backdrops behind each column. Defined before the arrow so we
  % can position the arrow with equal gaps relative to the actual box edges.
  \begin{scope}[on background layer]
    \node[fill=myblue!15, draw=myblue!60, rounded corners=2pt, inner sep=3mm,
          fit=(m-1-1)(m-2-1)(m-3-1)] (topL) {};
    \node[fill=mypurple!15, draw=mypurple!60, rounded corners=2pt, inner sep=3mm,
          fit=(m-1-2)(m-2-2)(m-3-2)
              (m-4-2.west|-m-2-2)(m-4-2.east|-m-2-2)] (topR) {};
    % Row-4 boxes: project each top-row cell's west and east onto row 4's y
    % so the fit box spans the full horizontal extent of the column above,
    % matching the width of the corresponding top box exactly.
    \node[fill=myblue!15, draw=myblue!60, rounded corners=2pt, inner sep=3mm,
          fit=(m-4-1)
              (m-1-1.west|-m-4-1)(m-1-1.east|-m-4-1)
              (m-2-1.east|-m-4-1)(m-3-1.east|-m-4-1)] (botL) {};
    \node[fill=mypurple!15, draw=mypurple!60, rounded corners=2pt, inner sep=3mm,
          fit=(m-4-2)
              (m-1-2.west|-m-4-2)(m-1-2.east|-m-4-2)
              (m-2-2.east|-m-4-2)(m-3-2.east|-m-4-2)] (botR) {};
  \end{scope}

  % Arrow sits between the two top boxes with equal left/right gaps, and
  % is vertically centered between the top and bottom blocks.
  \path let \p1=(m-1-1.north), \p2=(m-4-1.south),
            \p3=(topL.east), \p4=(topR.west) in
    coordinate (arrL) at ({\x3+5mm}, {(\y1+\y2)/2})
    coordinate (arrR) at ({\x4-5mm}, {(\y1+\y2)/2});
  \draw[->, thin] (arrL) -- (arrR)
    node[midway, above=1pt] (plab)
      {$p_0 := \mathtt{(+\ X\ (*\ Y\ Y))}$}
    node[midway, below=1pt] (flab)
      {$f_0 = \lambda\,\mathtt{x\ y.\ (+\ x\ (*\ y\ y))}$};
  % Size annotation for the abstraction body itself.
  \node[circle, fill=orange!40!white, draw=none,
        inner sep=0.6pt, font=\scriptsize]
    at ($(plab.north east)+(1mm,-0.5mm)$) {5};

  % Panel labels in bold, roughly top-aligned with each block, to its left.
  \node[anchor=north east, font=\bfseries]
    at ($(topL.north west)+(-2mm,-1mm)$) {(a)};
  \node[anchor=north east, font=\bfseries]
    at ($(botL.north west)+(-2mm,-1mm)$) {(b)};

  % AST-size annotations as small disks at the top-right corner of each cell.
  \foreach \r/\c/\s in {%
    1/1/9, 1/2/5,
    2/1/9, 2/2/5,
    3/1/9, 3/2/6,
    4/1/7, 4/2/7%
  } {
    \node[circle, fill=orange!40!white, draw=none,
          inner sep=0.6pt, font=\scriptsize]
      at (m-\r-\c.north east) {\s};
  }

\end{tikzpicture}
